% GENERATED FILE. Edit canonical lessons, then run npm run generate:books.
% canonical-source-hash: fnv1a64:a63766f599aaf348
% canonical-lessons: IT-C520-recitare, IT-C520-copione, IT-C520-ruolo, IT-C520-debuttare, IT-C520-regista

\chapter{Sul palco}
\label{ch:it-sul-palco}
\hlchaptermodality{520}

\begin{chapteropening}
\textbf{By the end of this chapter:} \emph{I can talk about acting in a play.}

Complete the last lesson of chapter 520: I can talk about acting in a play.
\end{chapteropening}

\section[recitare]{recitare — to act; to recite}

\label{lesson:IT-C520-recitare}



\begin{quote}
Before the new one: say the Italian for gallery, then the Italian for applause.
\end{quote}

\subsection*{You'll want to know: recitare}
\textbf{recitare} — \textquotedblleft{}to act; to recite\textquotedblright{}.

\begin{quote}
Recito in una compagnia di teatro amatoriale.
\end{quote}

\begin{quote}
\emph{I act in an amateur theatre company.}
\end{quote}

More than to recite: on stage it means to act.

\begin{cousinweb}
\emph{recitare} is a learned borrowing from Latin \emph{recitāre} \textquotedblleft{}to read aloud\textquotedblright{} (EtymDB), from \emph{citāre} \textquotedblleft{}to call\textquotedblright{}.
\end{cousinweb}

\begin{grammarlens}[title={What you've built}]
One more word for telling, explaining and describing at B1.
\end{grammarlens}

\subsection*{Your turn}
\begin{itemize}
\raggedright
  \item \emph{Say it:} \emph{recitare}
  \item \emph{Say it:} \emph{recitare}, once more
  \item \emph{Say it:} a sentence of your own with \emph{recitare}
  \item \emph{Recall:} say the Italian for gallery, then the Italian for applause, then say \emph{recitare} again
\end{itemize}

\begin{tcolorbox}[breakable,colback=teal!4,colframe=teal!35!black,title={Before you move on}]
What does recitare mean? (\textquotedblleft{}To act; to recite\textquotedblright{}.) Say it once more.
\end{tcolorbox}
\section[il copione]{il copione — script}

\label{lesson:IT-C520-copione}



\begin{quote}
Before the new one: say the Italian for applause, then the Italian for to act.
\end{quote}

\subsection*{You'll want to know: il copione}
\textbf{il copione} — \textquotedblleft{}script\textquotedblright{}.

\begin{quote}
Devo imparare il copione a memoria.
\end{quote}

\begin{quote}
\emph{I have to learn the script by heart.}
\end{quote}

Also a set pattern: \emph{seguire il copione}.

\begin{cousinweb}
\emph{copione} is built on \emph{copia} \textquotedblleft{}copy\textquotedblright{} (EtymDB), from Latin \emph{cōpia} \textquotedblleft{}plenty\textquotedblright{}: the big copy each actor works from.
\end{cousinweb}

\begin{grammarlens}[title={What you've built}]
One more word for telling, explaining and describing at B1.
\end{grammarlens}

\subsection*{Your turn}
\begin{itemize}
\raggedright
  \item \emph{Say it:} \emph{il copione}
  \item \emph{Say it:} \emph{il copione}, once more
  \item \emph{Say it:} a sentence of your own with \emph{il copione}
  \item \emph{Recall:} say the Italian for applause, then the Italian for to act, then say \emph{il copione} again
\end{itemize}

\begin{tcolorbox}[breakable,colback=teal!4,colframe=teal!35!black,title={Before you move on}]
What does il copione mean? (\textquotedblleft{}Script\textquotedblright{}.) Say it once more.
\end{tcolorbox}
\section[il ruolo]{il ruolo — role, part}

\label{lesson:IT-C520-ruolo}



\begin{quote}
Before the new one: say the Italian for to act, then the Italian for script.
\end{quote}

\subsection*{You'll want to know: il ruolo}
\textbf{il ruolo} — \textquotedblleft{}role, part\textquotedblright{}.

\begin{quote}
Nello spettacolo ho il ruolo del padre.
\end{quote}

\begin{quote}
\emph{In the play I have the part of the father.}
\end{quote}

Also at work: \emph{il mio ruolo in azienda}.

\begin{cousinweb}
\emph{ruolo} was borrowed from French \emph{rôle} \textquotedblleft{}a roll, a list\textquotedblright{}, from Medieval Latin \emph{rotulus} \textquotedblleft{}a roll of parchment\textquotedblright{}, from \emph{rota} \textquotedblleft{}wheel\textquotedblright{} (Lewis \& Short).
\end{cousinweb}

\begin{grammarlens}[title={What you've built}]
One more word for telling, explaining and describing at B1.
\end{grammarlens}

\subsection*{Your turn}
\begin{itemize}
\raggedright
  \item \emph{Say it:} \emph{il ruolo}
  \item \emph{Say it:} \emph{il ruolo}, once more
  \item \emph{Say it:} a sentence of your own with \emph{il ruolo}
  \item \emph{Recall:} say the Italian for to act, then the Italian for script, then say \emph{il ruolo} again
\end{itemize}

\begin{tcolorbox}[breakable,colback=teal!4,colframe=teal!35!black,title={Before you move on}]
What does il ruolo mean? (\textquotedblleft{}Role, part\textquotedblright{}.) Say it once more.
\end{tcolorbox}
\section[debuttare]{debuttare — to make one's debut}

\label{lesson:IT-C520-debuttare}



\begin{quote}
Before the new one: say the Italian for script, then the Italian for role.
\end{quote}

\subsection*{You'll want to know: debuttare}
\textbf{debuttare} — \textquotedblleft{}to make one's debut\textquotedblright{}.

\begin{quote}
Ha debuttato a teatro a vent'anni.
\end{quote}

\begin{quote}
\emph{She made her stage debut at twenty.}
\end{quote}

The noun is \emph{il debutto}.

\begin{cousinweb}
\emph{debuttare} was borrowed from French \emph{débuter} \textquotedblleft{}to make the first move\textquotedblright{} (EtymDB), built on \emph{but} \textquotedblleft{}a target\textquotedblright{}.
\end{cousinweb}

\begin{grammarlens}[title={What you've built}]
One more word for telling, explaining and describing at B1.
\end{grammarlens}

\subsection*{Your turn}
\begin{itemize}
\raggedright
  \item \emph{Say it:} \emph{debuttare}
  \item \emph{Say it:} \emph{debuttare}, once more
  \item \emph{Say it:} a sentence of your own with \emph{debuttare}
  \item \emph{Recall:} say the Italian for script, then the Italian for role, then say \emph{debuttare} again
\end{itemize}

\begin{tcolorbox}[breakable,colback=teal!4,colframe=teal!35!black,title={Before you move on}]
What does debuttare mean? (\textquotedblleft{}To make one's debut\textquotedblright{}.) Say it once more.
\end{tcolorbox}
\section[il regista]{il regista — director}

\label{lesson:IT-C520-regista}



\begin{quote}
Before the new one: say the Italian for role, then the Italian for to make one's debut.
\end{quote}

\subsection*{You'll want to know: il regista}
\textbf{il regista} — \textquotedblleft{}director\textquotedblright{}.

\begin{quote}
Il regista ha cambiato la fine dello spettacolo.
\end{quote}

\begin{quote}
\emph{The director changed the ending of the play.}
\end{quote}

The same form for a woman: \emph{la regista}. The work is \emph{la regia}.

\begin{cousinweb}
\emph{regista} is built on \emph{regia} \textquotedblleft{}direction\textquotedblright{}, usually traced to French \emph{régie} \textquotedblleft{}management\textquotedblright{}, from \emph{régir} \textquotedblleft{}to govern\textquotedblright{}, from Latin \emph{regere} \textquotedblleft{}to rule\textquotedblright{} (EtymDB, under French \emph{régie}).
\end{cousinweb}

\begin{grammarlens}[title={What you've built}]
One more word for telling, explaining and describing at B1.
\end{grammarlens}

\subsection*{Your turn}
\begin{itemize}
\raggedright
  \item \emph{Say it:} \emph{il regista}
  \item \emph{Say it:} \emph{il regista}, once more
  \item \emph{Say it:} a sentence of your own with \emph{il regista}
  \item \emph{Recall:} say the Italian for role, then the Italian for to make one's debut, then say \emph{il regista} again
\end{itemize}

\begin{tcolorbox}[breakable,colback=teal!4,colframe=teal!35!black,title={Before you move on}]
What does il regista mean? (\textquotedblleft{}Director\textquotedblright{}.) Say it once more.
\end{tcolorbox}
